\begingroup
\renewcommand{\arraystretch}{1.45}
\Needspace{22\baselineskip}
\InterpolationArticle{26}
If the method of Articles 20, 22 had been applied directly to all $\pi$ proposed values, without first dividing them into $\nu$ smaller periods, a function ($X'''$) of the $\tfrac12\pi-\tfrac12^{\mathrm{th}}$ or $\tfrac12\pi^{\mathrm{th}}$ order would have resulted, according as $\pi$ is odd or even; we shall now compare it with the function $X''$. Since the latter, when duly reduced, evidently reaches the $\mu n+m^{\mathrm{th}}$ order, we shall distinguish three cases.

I. When both $\mu$ and $\nu$ are odd, $\mu n+m=\tfrac12\mu\nu-\tfrac12=\tfrac12\pi-\tfrac12$ will hold. Hence it is clear that $X''$ must be entirely identical with $X'''$.

\Needspace{31\baselineskip}
II. When $\mu$ is even and $\nu$ odd, $\mu n+m=\tfrac12\mu\nu=\tfrac12\pi$, and therefore $X''$ is of the sa\SourceLeaf{314}%
me order as $X'''$. It therefore suffices to compare the last terms of the two functions. It is not difficult to see that these last terms in $X''$ arise solely from the terms $\gamma^m\cos mx+\delta^m\sin mx$, namely from those terms of the coefficients $\gamma^m,\delta^m$ that contain $\cos ny$ and $\sin ny$, or $\cos\mu nx$ and $\sin\mu nx$. Let these terms be
\[
\begin{aligned}
&\text{in }\gamma^m\quad.\,.\,.\,.\quad k\cos\mu nx+l\sin\mu nx\\[5pt]
&\text{in }\delta^m\quad.\,.\,.\,.\quad k'\cos\mu nx+l'\sin\mu nx
\end{aligned}
\]
from which there results in $X''$
\[
\begin{aligned}
&(k\cos\mu nx+l\sin\mu nx)\cos mx\\[5pt]
&\qquad{}+(k'\cos\mu nx+l'\sin\mu nx)\sin mx
\end{aligned}
\]
This gives the following terms of the $\mu n+m^{\mathrm{th}}$ order:
\[
\begin{aligned}
&\tfrac12(k-l')\cos(\mu n+m)x+\tfrac12(l+k')\sin(\mu n+m)x
\end{aligned}
\]
Now it can indeed be shown without difficulty \textit{from the formation of the coefficients} that
\[
\begin{aligned}
&(k-l')\sin(\mu n+m)a=(l+k')\cos(\mu n+m)a
\end{aligned}
\]
\Needspace{42\baselineskip}
but we prefer the following method, which seems more elegant. From Article 24 it follows that
\[
\begin{aligned}
&\gamma^m\sin\tfrac12y-\delta^m\cos\tfrac12y=0
\end{aligned}
\]
\textit{both} when for $y$ the value $\mu a$ is substituted, and for $\gamma^m,\delta^m$ the values derived from the first period,
\textit{and} when for $y$ the value $\mu a'$ is substituted, and for $\gamma^m,\delta^m$ the values derived from the second period, etc. Thus, if for $\gamma^m,\delta^m$ the undetermined expressions that have been found are substituted, so that $\gamma^m\sin\tfrac12y-\delta^m\cos\tfrac12y$ becomes a function of the angle $y$ of the form $G$ reaching the $n+\tfrac12=\tfrac12\nu$ order, this function will, for all $\nu$ values of $y$, namely $\mu a,\mu a',\mu a''$ etc., become $=0$, and will therefore, by Article 14, be divisible by
\[
\begin{aligned}
&\sin\tfrac12(y-\mu a)\sin\tfrac12(y-\mu a')\sin\tfrac12(y-\mu a'')\text{ etc.}
\end{aligned}
\]
and hence also by $\sin\tfrac12\nu(y-\mu a)$, and will consequently be of the form $h\sin\tfrac12\nu(y-\mu a)$, where $h$ is a fixed quantity. Thus the last terms of the expanded function $\gamma^m\sin\tfrac12y-\delta^m\cos\tfrac12y$ must be
\[
\begin{aligned}
&-h\sin\tfrac12\pi a\cos\tfrac12\nu y+h\cos\tfrac12\pi a\sin\tfrac12\nu y
\end{aligned}
\]
and all preceding terms must vanish; but these last terms clearly become
\SourceLeaf{315}
\[
\begin{aligned}
&=-\tfrac12(l+k')\cos(n+\tfrac12)y+\tfrac12(k-l')\sin(n+\tfrac12)y
\end{aligned}
\]
Consequently, it will follow that
\[
\begin{aligned}
&\tfrac12(k-l')=h\cos\tfrac12\pi a=h\cos\tfrac12(\mu n+m)a\\[5pt]
&\tfrac12(l+k')=h\sin\tfrac12\pi a=h\sin\tfrac12(\mu n+m)a\qquad\text{Q.~E.~D.}
\end{aligned}
\]
This equation, however, will be precisely the criterion for the last coefficients of the function
$X'''$, as shown in Article 24; hence in this case too the function $X''$ will be
entirely identical with $X'''$.

\Needspace{32\baselineskip}
III. When $\nu$ is even, $\mu n+m=\tfrac12\pi+m$, and therefore $X''$ is of higher order
than $X'''$. It is easy to see, however, that all terms in $X''$ exceeding the
$\tfrac12\pi^{\mathrm{th}}$ order arise solely from those terms of the individual
coefficients $\gamma',\delta',\gamma'',\delta''$ etc. that contain $\cos ny$ and
$\sin ny$, or $\cos\mu nx$ and $\sin\mu nx$. In the expression for
any coefficient $\gamma^\lambda$ or $\delta^\lambda$, let these terms be $k\cos\mu nx+l\sin\mu nx$;
from them there results in $X''$
\[
\begin{aligned}
&(k\cos\mu nx+l\sin\mu nx)\cos\lambda x\quad\text{in the first case}\\[5pt]
&\text{or}\quad(k\cos\mu nx+l\sin\mu nx)\sin\lambda x\quad\text{in the second case}
\end{aligned}
\]
and hence
\[
\begin{aligned}
&\tfrac12k\cos(\mu n+\lambda)x+\tfrac12l\sin(\mu n+\lambda)x\\[5pt]
&+\tfrac12k\cos(\mu n-\lambda)x+\tfrac12l\sin(\mu n-\lambda)x
\end{aligned}
\]
in the first case, and
\[
\begin{aligned}
&-\tfrac12l\cos(\mu n+\lambda)x+\tfrac12k\sin(\mu n+\lambda)x\\[5pt]
&+\tfrac12l\cos(\mu n-\lambda)x-\tfrac12k\sin(\mu n-\lambda)x
\end{aligned}
\]
in the second case. By Article 24, however, for $\cos(\mu n+\lambda)x$,
or $\cos(\tfrac12\pi+\lambda)x$, we must substitute $\cos\pi a\cos(\tfrac12\pi-\lambda)x+\sin\pi a\sin(\tfrac12\pi-\lambda)x$;
and for $\sin(\mu n+\lambda)x$, or $\sin(\tfrac12\pi+\lambda)x$,
$\sin\pi a\cos(\tfrac12\pi-\lambda)x-\cos\pi a\sin(\tfrac12\pi-\lambda)x$.
We shall therefore have in the first case
\[
\begin{aligned}
&\tfrac12(k+k\cos\pi a+l\sin\pi a)\cos(\tfrac12\pi-\lambda)x\\[5pt]
&+\tfrac12(l+k\sin\pi a-l\cos\pi a)\sin(\tfrac12\pi-\lambda)x
\end{aligned}
\]
and in the second case
\[
\begin{aligned}
&\tfrac12(l-l\cos\pi a+k\sin\pi a)\cos(\tfrac12\pi-\lambda)x\\[5pt]
&-\tfrac12(k+l\sin\pi a+k\cos\pi a)\sin(\tfrac12\pi-\lambda)x
\end{aligned}
\]
\SourceLeaf{316}
But from the equation $k\sin\tfrac12\pi a=l\cos\tfrac12\pi a$, which must hold in both cases
between $k$ and $l$ (Article 24), it follows without difficulty that
\[
\begin{aligned}
&\tfrac12(k+k\cos\pi a+l\sin\pi a)=k\\[5pt]
&\tfrac12(l+k\sin\pi a-l\cos\pi a)=l
\end{aligned}
\]
\Needspace{32\baselineskip}
From this it finally follows that, in place of those parts of $X''$ which would arise from the last terms
\[
\begin{aligned}
&k\cos ny+l\sin ny
\end{aligned}
\]
of any coefficient $\gamma^\lambda$ or $\delta^\lambda$,
with $\lambda$ denoting an index greater than 0, the following may be directly substituted
\[
\begin{aligned}
&k\cos(\tfrac12\pi-\lambda)x+l\sin(\tfrac12\pi-\lambda)x\\[5pt]
&\qquad\text{in the former case, or}\\[5pt]
&l\cos(\tfrac12\pi-\lambda)x-k\sin(\tfrac12\pi-\lambda)x\\[5pt]
&\qquad\text{in the latter case}
\end{aligned}
\]
In this way, instead of the function $X''$, which would reach the $\tfrac12\pi+m^{\mathrm{th}}$
order, we obtain another, $X''''$, which does not exceed the $\tfrac12\pi^{\mathrm{th}}$
order. It therefore remains only to compare the coefficients of the cosine and sine of the angle
$\tfrac12\pi x$, which have evidently remained the same in this function as in $X''$, with
the corresponding coefficients in $X'''$. A little attention shows, however,
that these terms arise solely from the first coefficient $\gamma$, namely
from its last terms; if these are represented by $k\cos ny+l\sin ny$, the former terms will be
$k\cos\tfrac12\pi x+l\sin\tfrac12\pi x$, and hence necessarily identical with
the corresponding terms in $X'''$. Thus it follows that the function $X''''$, obtained by
the rules just given in place of the function $X''$, becomes entirely identical with the function $X'''$.

\endgroup
